% !TEX root = task-2-10-8.tex
% Editable task bank. Append new BankTask blocks at the end of the document.
\documentclass[12pt,a4paper]{article}
\usepackage{physics-textbook}
\usepackage{graphicx,needspace}
\usetikzlibrary{decorations.pathmorphing}
\newif\ifShowSolutions
\ifdefined\TutorEdition
  \ShowSolutionstrue
\else
  \ShowSolutionsfalse % Change to \ShowSolutionstrue to display solutions.
\fi
\newcounter{banktask}
\newsavebox{\BankTaskBox}
\newcommand{\BankTaskContent}[3]{%
  \noindent\textbf{\thebanktask.}\quad #2\par
  \ifShowSolutions
    \Needspace{50mm}% Keep the solution heading with its opening text or diagram.
    #3\par
    \ifdefined\TutorEdition
      \ifcsname TutorSource#1\endcsname
        \nopagebreak\smallskip{\small\csname TutorSource#1\endcsname\par}%
      \fi
    \fi
  \fi
}
\newcommand{\BankTask}[3]{%
  \refstepcounter{banktask}\par\medskip
  \sbox{\BankTaskBox}{\begin{minipage}{\linewidth}%
    \BankTaskContent{#1}{#2}{#3}\end{minipage}}%
  \ifdim\dimexpr\ht\BankTaskBox+\dp\BankTaskBox\relax>\dimexpr\textheight-20mm\relax
    \Needspace{35mm}\BankTaskContent{#1}{#2}{#3}%
  \else
    \Needspace{\dimexpr\ht\BankTaskBox+\dp\BankTaskBox\relax}%
    \noindent\usebox{\BankTaskBox}\par
  \fi
  \medskip
}
\begin{document}
\BookSubsection{2.10.8}{График конденсации}


% Source E2B7DC: https://ege.fipi.ru/bank/questions.php?search=1&proj=BA1F39653304A5B041B656915DC36B38&qid=E2B7DC
\BankTask{E2B7DC}{
В начальный момент в сосуде под лёгким поршнем находится только жидкий эфир. На рисунке схематично представлен график зависимости температуры \(t\) эфира от времени \(\tau\) его нагревания и последующего охлаждения.
\begin{center}
\begin{tikzpicture}[x=0.85cm,y=0.65cm,>=stealth]
\draw[->] (0,0)--(6.8,0) node[below] {$\tau$};
\draw[->] (0,0)--(0,4.7) node[left] {$t,\,{}^\circ\mathrm{C}$};
\node[below left] at (0,0) {$0$};
\foreach \y/\v in {1/20,2/40,3/60}{\draw (-0.08,\y)--(0.08,\y);\node[left] at (-0.1,\y) {$\v$};}
\foreach \y in {0.5,1.5,2.5}{\draw (-0.06,\y)--(0.06,\y);}
\draw[dashed] (0,2)--(3.35,2);
\draw (0,1)--(0.75,2)--(2.1,2)--(2.72,3.4)--(3.35,2)--(4.7,2)--(5.35,1);
\node[below right] at (0,1) {$A$};
\node[above] at (0.75,2) {$B$};
\node[above left] at (2.1,2) {$C$};
\node[above] at (2.72,3.4) {$D$};
\node[above right] at (3.35,2) {$E$};
\node[above right] at (4.7,2) {$F$};
\node[right] at (5.35,1) {$G$};
\end{tikzpicture}
\end{center}
\begin{enumerate}
\item Температура кипения эфира равна \(60\,{}^\circ\mathrm{C}\).
\item В точке \(F\) в сосуде находился только жидкий эфир.
\item На участке \(BC\) внутренняя энергия эфира увеличивалась.
\item В точке \(G\) эфир отвердел.
\item Время, за которое весь эфир выкипел, меньше, чем время, за которое он сконденсировался.
\end{enumerate}
}{
\smallskip\noindent\textbf{Решение от GPT}\par
1) Неверно: горизонтальный участок $BC$ соответствует кипению при $40\,{}^\circ\text{C}$.

2) Верно: в точке $F$ конденсация завершилась, далее охлаждается жидкость.

3) Верно: при кипении внутренняя энергия увеличивается за счёт увеличения потенциальной энергии взаимодействия молекул.

4) Неверно: участок $FG$ соответствует охлаждению жидкости; признаков отвердевания нет.

5) Неверно: участки $BC$ и $EF$ на схематическом графике имеют одинаковую длительность.

Ответ: $23$.
}

\end{document}
